% sec-01: the AMP model.  Sources: chunk-04 (FNIH and AMP held up as the model,
% AMP-AD validated 20 targets), chunk-03 ($200B), chunk-08 (cost share).
\section{The Partnership Model the Report Names}

\textit{Science: A New Golden Age} holds up the Foundation for the NIH and the
Accelerating Medicines Partnership as its model for public-private structure,
noting that the Alzheimer's partnership, one of twelve disease-focused AMPs,
experimentally validated twenty candidate drug targets
\cite{kratsios2026sciencegoldenage}. In the same chapter the report observes
that American discoveries are increasingly tested abroad because per-patient
trial costs run far higher here, and its annex asks agencies to prioritize
non-federal cost share so that federal dollars draw in private investment rather
than standing alone \cite{kratsiosvought2026fy2028priorities}.

This concept applies that structure to a narrow, pre-competitive question in
resectable pancreatic cancer: what happens pharmacologically when a RAS(ON)
multi-selective inhibitor is given around a curative-intent resection. No single
party can answer it. The drug developer holds the compound, an academic center
holds the patients and the operative expertise, and the question that matters,
what the resection specimen shows, is a public good that no company will fund
on its own because it does not differentiate a product.

\begin{apptable}
\begin{tabularx}{\textwidth}{>{\raggedright\arraybackslash}p{4.05cm} >{\raggedright\arraybackslash}X}
\headrow \thead{AMP feature} & \thead{Instance proposed here}\\
\midrule
A pre-competitive question no party owns & Perioperative RAS pathway pharmacodynamics measured in resected human PDAC tissue \\
Shared governance with public reporting & Protocol, analysis plan, and negative results deposited before database lock \\
Non-federal cost share drawing in private investment & Investigational drug supply and site infrastructure contributed rather than purchased \\
Data released to the field, not to a sponsor & Specimen-level pharmacodynamic dataset released under an open license \\
\end{tabularx}
\end{apptable}
\tabcap{Four defining features of the Accelerating Medicines Partnership model,
and the specific instance each takes in a perioperative RAS consortium built
around a resectable pancreatic cancer trial.}
